% ch24-community-block
\section{المجتمع والحظر: كيف يعمل بالكامل}

مجتمع \en{ScholarPath} هو المساحة التي يتناقش فيها الطلاب: يسألون عن المنح، يشاركون
تجاربهم، ويصوّتون على أفضل الإجابات. لكن أيّ مساحة مفتوحة تحتاج قواعد تحكمها: مَن
يكتب؟ كيف نمنع المحتوى الضارّ؟ وماذا يحدث حين يريد طالبٌ ألّا يرى شخصًا آخر إطلاقًا؟

في هذا الفصل هنفهم ثلاثة أشياء متتالية. أولًا: كيانات المجتمع — الفئات والمنشورات
والردود والتصويت والإبلاغ والإشارات المرجعيّة — وكيف تترابط. ثانيًا: دورة حياة منشور
كاملة، من لحظة كتابته حتى إخفائه تلقائيًّا لو أبلغ عنه عدد كافٍ من الناس. ثالثًا: نظام
الحظر (\en{Block}) — وهو سؤالنا الأهمّ: كيف يُنشَأ الحظر، والأهمّ: أين يُنفَّذ؟
هنكتشف أنّ الحظر لا يُطبَّق في مكان واحد، بل دفاعيًّا في كل مسار: فحص عند الكتابة،
وفلتر عند القراءة.

نفترض أنّك تعرف من الفصول السابقة فكرة \en{CQRS} (كل عملية إمّا أمر أو استعلام)،
و\en{soft delete} (الحذف المنطقيّ بعلَم بدل المحو الفعليّ)، و\en{SignalR} وقنوات البثّ
الحيّ (\en{hubs} ومجموعات المشتركين)، وتنقية \en{HTML} قبل التخزين. سنشير إليها بسطر
عند الحاجة فقط.

% ====================================================================
\subsection{كيانات المجتمع: الخريطة الكاملة}

المجتمع مبنيّ على حفنة كيانات بسيطة. أذكى قرار في التصميم أنّ المنشور والردّ نفس
الكيان: كلاهما \code{ForumPost}. الفرق الوحيد حقلٌ اسمه \code{ParentPostId}: لو كان
\code{null} فهذا منشور جذر (\en{thread})، ولو كان يحمل مُعرّفًا فهذا ردّ على منشور آخر.
هذا يوفّر علينا جدولًا ومنطقًا مكرّرًا.

\begin{center}
\begin{tabular}{L{3.2cm} L{9cm}}
\toprule
الكيان & الدور \\
\midrule
\en{ForumCategory} & فئة نقاش (منح، تأشيرات، سكن...) لها \code{Slug} وعلَم \code{IsActive}. \\
\en{ForumPost} & منشور جذر أو ردّ (يفرّقهما \code{ParentPostId}). يحمل عدّادات مخزّنة و\code{ModerationStatus}. \\
\en{ForumVote} & صوت واحد لكل مستخدم على منشور: \code{Up} أو \code{Down}. \\
\en{ForumFlag} & بلاغ من مستخدم عن منشور، بسبب و\code{IsValid} يقدر الأدمن يبطله. \\
\en{ForumBookmark} & إشارة مرجعيّة تربط طالبًا بمنشور حفظه ليرجع إليه. \\
\bottomrule
\end{tabular}
\end{center}

المقتطف التالي من ملف \code{Community.cs} يبيّن قلب \code{ForumPost} — لاحظ حقلين مصيريّين:
\code{ModerationStatus} (حالة الإشراف) و\code{IsAutoHidden} (الإخفاء التلقائيّ):

\begin{english}
\begin{lstlisting}[language={[Sharp]C}]
public class ForumPost : AuditableEntity, ISoftDeletable
{
    public Guid AuthorId { get; set; }
    public Guid? CategoryId { get; set; }
    public Guid? ParentPostId { get; set; } // null = root thread; else a reply

    public PostModerationStatus ModerationStatus { get; set; }
        = PostModerationStatus.Visible;

    // Cached aggregates (kept in sync on write, cheap to read)
    public int UpvoteCount { get; set; }
    public int FlagCount { get; set; }
    public int ReplyCount { get; set; }

    public bool IsAutoHidden { get; set; } // hidden by the threshold rule
    public bool IsDeleted { get; set; }    // soft delete
    // ... navigations: Replies, Votes, Flags, Bookmarks ...
}
\end{lstlisting}
\end{english}

شرح سريع لأهمّ الأسطر:
\begin{itemize}
\item \code{ParentPostId} هو المفتاح الذكيّ: منشور وردّ في نفس الجدول، والفرق حقل واحد.
\item \code{ModerationStatus} حالة إشراف من نوع \code{PostModerationStatus}: \code{Visible}
  أو \code{Removed} (حذفه الأدمن) أو \code{PendingReview} (بانتظار مراجعة) أو \code{Hidden}.
\item \code{IsAutoHidden} علَمٌ منفصل: صحيح حين تتخطّى البلاغاتُ العتبةَ فيُخفى المنشور
  تلقائيًّا قبل أن يتدخّل أيّ إنسان.
\item العدّادات \code{UpvoteCount} و\code{FlagCount} و\code{ReplyCount} مخزّنة مسبقًا
  (\en{cached aggregates}): نحدّثها لحظة الكتابة حتى تكون القراءة رخيصة بلا عدٍّ متكرّر.
\item \code{IsDeleted} هو الـ \en{soft delete}: المنشور المحذوف يبقى صفًّا في القاعدة لكن
  يختفي من كل الاستعلامات.
\end{itemize}

الرسم التالي يجمع العلاقات: الفئة تحوي منشورات، والمنشور له ردود، وحوله التصويت والإبلاغ
والإشارة المرجعيّة.

\begin{center}
\begin{LTR}
\begin{tikzpicture}[node distance=0.6cm and 1.4cm]
  \node[dbbox] (cat) {Category};
  \node[dbbox,right=1.7cm of cat] (post) {Post\\(root)};
  \node[dbbox,right=1.7cm of post] (reply) {Reply};
  \node[sbox,text width=2.0cm,above=0.95cm of post] (vote) {Vote\\Up / Down};
  \node[sbox,text width=2.0cm,below=0.95cm of post] (flag) {Flag};
  \node[sbox,text width=2.0cm,above=0.95cm of reply] (book) {Bookmark};
  \draw[arr] (cat) -- (post) node[lbl,midway,above] {posts};
  \draw[arr] (post) -- (reply) node[lbl,midway,above] {ParentPostId};
  \draw[arrb] (post) -- (vote);
  \draw[arrb] (post) -- (flag);
  \draw[arrb] (post) -- (book);
\end{tikzpicture}
\end{LTR}
\end{center}

% ====================================================================
\subsection{دورة حياة منشور: من الكتابة إلى العرض}

المنشور يبدأ حياته بأمر \code{CreatePostCommand}. أوّل خطوة حارسة: مَن يكتب؟ فقط الطالب.
المستشارون والجهات لا ينتجون محتوى هنا. ثانيًا: كل نصٍّ يمرّ على مُنقِّي \en{HTML} قبل
التخزين، حتى لا يُحقَن سكربت خبيث. المقتطف من \code{CreatePostCommand} يبيّن الحارسَين
معًا، ثم إطلاق حدث البثّ الحيّ:

\begin{english}
\begin{lstlisting}[language={[Sharp]C}]
if (!currentUser.IsInRole("Student"))
    throw new ForbiddenAccessException("Only students can create community posts.");

var sanitizer = new Ganss.Xss.HtmlSanitizer();
var titleEn = sanitizer.Sanitize(request.TitleEn);
var bodyEn  = sanitizer.Sanitize(request.BodyEn);
// ... TitleAr / BodyAr sanitized the same way (posts are bilingual) ...

var post = new ForumPost { AuthorId = authorId, CategoryId = request.CategoryId,
    TitleEn = titleEn, BodyEn = bodyEn /* ... */ };
db.ForumPosts.Add(post);

// Raise BEFORE save so the dispatcher (runs after SaveChanges) broadcasts it.
post.RaiseDomainEvent(new ForumPostCreatedEvent(post.Id, authorId, request.CategoryId));
await db.SaveChangesAsync(ct);
\end{lstlisting}
\end{english}

الأسطر المهمّة:
\begin{itemize}
\item الفحص \code{IsInRole("Student")} يرفض أيّ دور آخر برمي \code{ForbiddenAccessException}.
\item \code{HtmlSanitizer} ينظّف العنوان والنصّ باللغتين — منشورات المجتمع ثنائيّة اللغة
  مثل المنح تمامًا.
\item \code{RaiseDomainEvent} يُطلَق قبل \code{SaveChangesAsync} عمدًا: الموزّع
  يلتقط الحدث بعد الحفظ فيبثّ الإشعار الحيّ. لو أطلقناه بعد الحفظ لضاع البثّ.
\end{itemize}

بعد أن يُعرَض المنشور، يتفاعل الطلاب معه: يردّون، ويصوّتون. التصويت أمرٌ لطيف اسمه
\code{ToggleVoteCommand}: نفس الزرّ يضيف الصوت أو يزيله أو يقلبه. لاحظ قاعدتين: لا تصوّت
على محتواك، وعدّاد المنشور يتحدّث في نفس اللحظة.

\begin{english}
\begin{lstlisting}[language={[Sharp]C}]
if (post.AuthorId == currentUser.UserId)
    throw new ConflictException("You can't vote on your own content.");

var existingVote = post.Votes.FirstOrDefault(v => v.UserId == currentUser.UserId);
if (existingVote != null && existingVote.VoteType == request.VoteType)
{
    db.ForumVotes.Remove(existingVote);            // same vote again -> toggle off
    if (request.VoteType == VoteType.Up) post.UpvoteCount--; else post.DownvoteCount--;
}
// ... else: swap direction, or add a brand-new vote ...
await db.SaveChangesAsync(ct);
\end{lstlisting}
\end{english}

الفكرة أنّ العدّاد المخزّن \code{UpvoteCount} يتغيّر مع كل عمليّة، فلا نضطرّ لعدّ الأصوات
عند كل قراءة. والقيد الفريد على (المنشور، المستخدم) في القاعدة يضمن صوتًا واحدًا لكل شخص
مهما تكرّرت الطلبات.

\begin{keyp}[نقطة مهمة: العدّادات المخزّنة]
\code{UpvoteCount} و\code{FlagCount} و\code{ReplyCount} كلّها أرقام محفوظة على صفّ المنشور،
نحدّثها لحظة الكتابة. النتيجة: صفحة المجتمع تقرأ عددًا جاهزًا بدل \code{COUNT} على آلاف
الأصوات في كل تحميل. الثمن أن نبقيها متّسقة في كل أمر — وهو ما نراه في التصويت والإبلاغ.
\end{keyp}

% ====================================================================
\subsection{الإبلاغ والإخفاء التلقائيّ: حماية بلا تدخّل بشريّ فوريّ}

هنا تكمن أذكى قاعدة في المجتمع. حين يرى طالبٌ منشورًا مسيئًا، يبلّغ عنه بأمر
\code{FlagPostCommand}. لكن بلاغًا واحدًا لا يكفي لإخفاء منشور — قد يكون كيديًّا. القاعدة:
لو أبلغ ثلاثة أشخاص مختلفين (بلاغات صالحة متمايزة) عن نفس المنشور، يُخفى تلقائيًّا
ويُحوَّل إلى \code{PendingReview} بانتظار قرار الأدمن.

المقتطف من \code{FlagPostCommand} هو قلب القاعدة:

\begin{english}
\begin{lstlisting}[language={[Sharp]C}]
// distinct valid flaggers already persisted, + our own new one
var distinctValidFlags = priorDistinctValidFlags + 1;
post.FlagCount++;

// Auto-hide rule: 3 distinct valid flags route the post to moderation.
if (distinctValidFlags >= 3 && !post.IsAutoHidden)
{
    post.IsAutoHidden = true;
    post.AutoHiddenAt = DateTimeOffset.UtcNow;
    post.ModerationStatus = PostModerationStatus.PendingReview;
    post.RaiseDomainEvent(new PostAutoHiddenEvent(post.Id, distinctValidFlags));
}
await db.SaveChangesAsync(ct);
\end{lstlisting}
\end{english}

ملاحظات دقيقة:
\begin{itemize}
\item نعدّ المُبلِّغين المتمايزين (\code{Distinct}) لا عدد البلاغات — حتى لا يخفي
  شخصٌ واحد منشورًا بتكرار البلاغ. والفحص المسبق على تكرار المُبلِّغ نفسه يمنع ذلك أصلًا.
\item عند بلوغ الثلاثة يُرفع علَم \code{IsAutoHidden} وتتحوّل الحالة إلى \code{PendingReview}،
  ويُطلَق حدث \code{PostAutoHiddenEvent} لإخطار الأدمن.
\item الشرط \code{!post.IsAutoHidden} يمنع الإخفاء وإطلاق الحدث مرّتين لو أُعيدت المحاولة.
\end{itemize}

الرسم يلخّص القرار: هل بلغت البلاغاتُ المتمايزة الصالحة العتبةَ؟

\begin{center}
\begin{LTR}
\begin{tikzpicture}[node distance=0.6cm and 1.1cm]
  \node[pbox] (flag) {FlagPost};
  \node[sbox,text width=2.2cm,right=1.1cm of flag] (count) {distinct valid flags};
  \node[dbox,right=1.1cm of count] (dec) {flags $\geq 3$ ?};
  \node[fbox,above right=0.35cm and 1.2cm of dec] (hide) {AutoHidden\\PendingReview};
  \node[tbox,below right=0.35cm and 1.2cm of dec] (stay) {stays Visible};
  \draw[arr] (flag) -- (count);
  \draw[arr] (count) -- (dec);
  \draw[arr] (dec) -- node[lbl,above] {yes} (hide);
  \draw[arr] (dec) -- node[lbl,below] {no} (stay);
\end{tikzpicture}
\end{LTR}
\end{center}

بعد الإخفاء التلقائيّ يأتي دور الأدمن. أمامه مساران: \code{DismissPostFlags} (يقرّر أنّ
المنشور سليم فيُبطل كل البلاغات ويعيده \code{Visible})، أو \code{DeletePost} (يحذفه
منطقيًّا). المقتطف من \code{DismissPostFlags} يبيّن الرجوع للحالة الطبيعيّة:

\begin{english}
\begin{lstlisting}[language={[Sharp]C}]
foreach (var flag in post.Flags) flag.IsValid = false; // invalidate every flag
post.FlagCount = 0;
post.IsAutoHidden = false;
post.AutoHiddenAt = null;
post.ModerationStatus = PostModerationStatus.Visible;
post.ModeratedByAdminId = adminId;
post.ModeratedAt = DateTimeOffset.UtcNow;
await db.SaveChangesAsync(ct);
\end{lstlisting}
\end{english}

\begin{spotlight}[ركّز هنا لإبهار اللجنة: حماية ذاتيّة بلا انتظار]
عتبة البلاغ الثلاثيّ تجعل المجتمع يحمي نفسه. المحتوى المسيء يختفي لحظة يتّفق ثلاثة
مستخدمين مستقلّين عليه — بلا انتظار أدمن يقظ على مدار الساعة. وفي الوقت نفسه، اشتراط
ثلاثة مُبلِّغين متمايزين يمنع إساءة استخدام الزرّ لإسكات رأي واحد. إنّها موازنة
هندسيّة أنيقة بين سرعة الحماية وعدالة القرار، والقرار النهائيّ يبقى للأدمن (إبقاء أو حذف).
\end{spotlight}

% ====================================================================
\subsection{البثّ الحيّ عبر CommunityHub وبوّابات الرؤية}

حين يُنشر ردّ جديد على منشور، نريد أن يظهر فورًا لكل مَن يشاهد ذلك الموضوع دون تحديث
الصفحة. هذا يتمّ عبر \code{CommunityHub} (قناة \en{SignalR}). العميل يشترك في مجموعتين:
مجموعة الفئة (\code{JoinCategory}) ليعرف بالمنشورات الجديدة، ومجموعة المنشور
(\code{JoinPost}) ليعرف بالردود الجديدة على موضوع بعينه.

لكن الاشتراك ليس مفتوحًا للجميع. هنا تعمل بوّابة الرؤية \en{SEC-02}: لا تشترك إلّا في فئة
نشطة فعلًا، ولا في منشور مرئيّ فعلًا — غير محذوف، وغير مخفيّ تلقائيًّا، وحالته
\code{Visible}. المقتطف من \code{CommunityHub} يبيّن حارس \code{JoinPost}:

\begin{english}
\begin{lstlisting}[language={[Sharp]C}]
public async Task JoinPost(string postId)
{
    if (!Guid.TryParse(postId, out var id))
        throw new HubException("Invalid post id.");

    var visible = await db.ForumPosts.AsNoTracking()
        .AnyAsync(p => p.Id == id && !p.IsDeleted && !p.IsAutoHidden
            && p.ModerationStatus == PostModerationStatus.Visible);
    if (!visible)
        throw new HubException("Post not found.");

    await Groups.AddToGroupAsync(Context.ConnectionId, $"forum-post:{postId}");
}
\end{lstlisting}
\end{english}

لماذا نفحص الرؤية قبل الاشتراك؟ لأنّ بدون هذا الحارس يقدر أيّ مستخدم مُصادَق أن ينضمّ
لمجموعة \code{forum-post} لمنشور محذوف أو مخفيّ، فيتنصّت على ردوده الحيّة. الحارس يغلق
الباب عند لحظة الانضمام.

بعد الاشتراك، حين يُنشأ ردّ، يُطلَق حدث \code{ForumReplyCreatedEvent}، ومعالِجه ينادي
\code{CommunityRealtimeNotifier} الذي يبثّ للمجموعة فقط:

\begin{english}
\begin{lstlisting}[language={[Sharp]C}]
public Task NotifyNewReplyAsync(Guid replyId, Guid parentPostId, CancellationToken ct) =>
    hub.Clients
       .Group($"forum-post:{parentPostId}")
       .SendAsync("ReplyCreated", new { replyId, parentPostId }, ct);
\end{lstlisting}
\end{english}

الرسم يوضّح الـ \en{fan-out}: ردٌّ واحد يصل لكل المشتركين في مجموعة المنشور دفعة واحدة.

\begin{center}
\begin{LTR}
\begin{tikzpicture}[node distance=0.5cm and 1.0cm]
  \node[actor] (author) {Student reply};
  \node[pbox,right=1.0cm of author] (cmd) {CreateReply};
  \node[gbox,right=1.0cm of cmd] (hub) {CommunityHub};
  \node[dbbox,text width=2.4cm,right=1.0cm of hub] (grp) {group\\forum-post:id};
  \node[sbox,text width=1.9cm,above right=0.5cm and 1.0cm of grp] (s1) {subscriber};
  \node[sbox,text width=1.9cm,right=1.1cm of grp] (s2) {subscriber};
  \node[sbox,text width=1.9cm,below right=0.5cm and 1.0cm of grp] (s3) {subscriber};
  \draw[arr] (author) -- (cmd);
  \draw[arr] (cmd) -- node[lbl,above] {ReplyCreated} (hub);
  \draw[arr] (hub) -- (grp);
  \draw[arr] (grp) -- (s1);
  \draw[arr] (grp) -- (s2);
  \draw[arr] (grp) -- (s3);
\end{tikzpicture}
\end{LTR}
\end{center}

% ====================================================================
\subsection{الحظر: كيف يُنشَأ وأين يُنفَّذ}

نصل الآن للسؤال الأهمّ. حين يريد طالبٌ ألّا يتعامل مع مستخدم آخر، يُصدر
\code{BlockUserCommand}. هذا الأمر بسيط: ينشئ صفًّا واحدًا في جدول \code{UserBlock} يربط
الحاجب بالمحجوب.

\begin{english}
\begin{lstlisting}[language={[Sharp]C}]
if (request.UserIdToBlock == currentUser.UserId)
    throw new ConflictException("You cannot block yourself.");
// ... return early if the block row already exists (idempotent) ...

var block = new UserBlock
{
    BlockerId = currentUser.UserId ?? throw new ForbiddenAccessException(),
    BlockedUserId = request.UserIdToBlock,
    Reason = request.Reason
};
db.UserBlocks.Add(block);
await db.SaveChangesAsync(ct);
\end{lstlisting}
\end{english}

لكنّ إنشاء الصفّ نصف القصّة. الصفّ مجرّد بيان بأنّ هناك حظرًا؛ السؤال الحقيقيّ:
مَن يقرأ هذا الصفّ ويطبّقه؟ الجواب المهمّ: التنفيذ موزّع دفاعيًّا على كل مسار، وينقسم
لنوعين — فحص عند الكتابة، وفلتر عند القراءة.

\begin{idea}[الفكرة: تنفيذ الحظر = كتابة + قراءة]
عند \en{الكتابة}: كل أمر يُنتج تفاعلًا (إرسال رسالة، كتابة ردّ) يفحص جدول \code{UserBlock}
قبل التنفيذ ويرفض لو كان هناك حظر. عند \en{القراءة}: كل استعلام يعرض محتوى شخص آخر
يمرّ على فلتر يُسقط محتوى المحجوبين. لو غاب أيّ طرف، تسرّب المحتوى — لذا نطبّقه في الطرفين.
\end{idea}

المسار الأوّل (كتابة) — أمر \code{SendMessageCommand} يفحص الحظر قبل حفظ الرسالة:

\begin{english}
\begin{lstlisting}[language={[Sharp]C}]
var isBlocked = await db.UserBlocks.AnyAsync(b =>
    (b.BlockerId == recipientId          && b.BlockedUserId == currentUser.UserId) ||
    (b.BlockerId == currentUser.UserId   && b.BlockedUserId == recipientId), ct);

if (isBlocked)
    throw new ConflictException("Cannot send message. User is blocked.");
\end{lstlisting}
\end{english}

لاحظ أنّ الفحص ثنائيّ الاتجاه: يرفض سواء حجبتُ أنا الطرف الآخر أو حجبني هو. نفس
المنطق في المجتمع: أمر \code{CreateReplyCommand} ينادي \code{AreBlockedAsync} فيرفض الردّ
لو كان بينهما حظر في أيّ اتجاه.

المسار الثاني (قراءة) — الفلتر المشترك \code{CommunityBlockFilter.NotBlockedWith}. هو تعبير
\code{Expression} واحد تستخدمه كل استعلامات المجتمع، حتى لا تنحرف دلالة الحظر بين مكان
وآخر:

\begin{english}
\begin{lstlisting}[language={[Sharp]C}]
public static Expression<Func<ForumPost, bool>> NotBlockedWith(
    IApplicationDbContext db, Guid viewerId) =>
    p => !db.UserBlocks.Any(b =>
        (b.BlockerId == viewerId    && b.BlockedUserId == p.AuthorId) ||
        (b.BlockerId == p.AuthorId  && b.BlockedUserId == viewerId));
\end{lstlisting}
\end{english}

هذا التعبير يعني: «أرِني المنشور فقط إن لم يكن صاحبه في علاقة حظر معي في أيّ اتّجاه».
تستدعيه \code{GetPostsQuery} على قائمة المنشورات، و\code{GetMyBookmarksQuery} يطبّق نفس
الشرط بيده فيُسقِط من إشاراتي المرجعيّة أيّ منشور صار صاحبه محجوبًا. الجدول التالي يجمع
نقاط التنفيذ:

\begin{center}
\begin{tabular}{L{3.4cm} L{2.6cm} L{6cm}}
\toprule
الموضع & النوع & السلوك \\
\midrule
\en{SendMessage} & كتابة & يرفض إرسال الرسالة. \\
\en{CreateReply} & كتابة & يرفض الردّ (\code{AreBlockedAsync}). \\
\en{GetPosts} & قراءة & يُخفي منشورات المحجوب من القائمة. \\
\en{GetMyBookmarks} & قراءة & يُسقط المحجوب من الإشارات المرجعيّة. \\
\en{GetConversations} & قراءة & يعلّم مَن حجبتُ لعرض زرّ إلغاء الحظر. \\
\bottomrule
\end{tabular}
\end{center}

الرسم يجسّد الفكرة: أمرٌ واحد ينشئ صفًّا واحدًا، ثم يُقرأ في مسارين متعامدين.

\begin{center}
\begin{LTR}
\begin{tikzpicture}[node distance=0.55cm and 1.1cm]
  \node[actor] (u) {Blocker};
  \node[pbox,right=1.1cm of u] (cmd) {BlockUser};
  \node[dbbox,text width=2.3cm,right=1.1cm of cmd] (row) {UserBlock row};
  \node[fbox,text width=2.9cm,above right=0.45cm and 1.3cm of row] (write) {write path\\SendMessage rejected};
  \node[tbox,text width=2.9cm,below right=0.45cm and 1.3cm of row] (read) {read path\\BlockFilter hides};
  \draw[arr] (u) -- (cmd);
  \draw[arr] (cmd) -- (row);
  \draw[arr] (row) -- node[lbl,above] {on write} (write);
  \draw[arr] (row) -- node[lbl,below] {on read} (read);
\end{tikzpicture}
\end{LTR}
\end{center}

\begin{secnote}
الفحص ثنائيّ الاتجاه مقصود ومهمّ. لو كان الفلتر يفحص «مَن حجبتُ أنا» فقط، لتمكّن مَن حجبني
من متابعة محتواي والردّ عليّ. توحيد الدلالة في تعبير \code{NotBlockedWith} واحد يضمن أنّ
كل استعلام يطبّق النسخة المتبادلة نفسها، فلا تنشأ ثغرة من اختلاف التطبيق بين شاشة وأخرى.
\end{secnote}

\begin{spotlight}[ركّز هنا لإبهار اللجنة: الحظر دفاعيّ في كل مسار]
لو سألتك اللجنة «أين يُنفَّذ الحظر؟» فالإجابة القويّة ليست «في مكان واحد»، بل: في كل
مسار كتابة وقراءة على حِدة. أمر واحد (\code{BlockUser}) يكتب صفًّا واحدًا، لكن هذا الصفّ
يُقرأ في \code{SendMessage} و\code{CreateReply} (كتابة) وفي \code{GetPosts} و\code{GetMyBookmarks}
و\code{GetConversations} (قراءة). هذا نمط \en{defense-in-depth}: لا نعتمد على طبقة واحدة
تحرس كل شيء، بل نغلق الباب عند كل نقطة تماسّ. النتيجة: لا مسار واحد يتسرّب منه محتوى المحجوب.
\end{spotlight}

\begin{qa}[سؤال متوقّع: لماذا لا تحرس مكانًا مركزيًّا واحدًا فقط؟]
لأنّ المسارات مختلفة الطبيعة: الكتابة تحتاج رفضًا صريحًا برسالة خطأ، والقراءة تحتاج
إسقاطًا صامتًا من النتائج. لا توجد نقطة مركزيّة واحدة يمرّ منها الاثنان. لذلك نضع
الفحص عند كل أمر كتابة، والفلتر (\code{NotBlockedWith}) عند كل استعلام قراءة — مع توحيد
الدلالة المتبادلة في مكان واحد حتى لا تتكرّر ولا تنحرف.
\end{qa}

% ====================================================================
\begin{summary}[الخلاصة في سطور]
\begin{itemize}
\item المنشور والردّ نفس الكيان \code{ForumPost}؛ يفرّقهما \code{ParentPostId}.
\item حول المنشور: تصويت (\code{ForumVote}) وبلاغ (\code{ForumFlag}) وإشارة مرجعيّة
  (\code{ForumBookmark})، وعدّادات مخزّنة تجعل القراءة رخيصة.
\item الإنشاء يحرسه دورٌ (طالب فقط) وتنقية \en{HTML}، ويُطلق حدثًا للبثّ الحيّ قبل الحفظ.
\item الإبلاغ يحمي المجتمع ذاتيًّا: ثلاثة مُبلِّغين متمايزين يُخفون المنشور تلقائيًّا
  (\code{IsAutoHidden}) ويحوّلونه \code{PendingReview}، والقرار النهائيّ للأدمن.
\item البثّ الحيّ عبر \code{CommunityHub} مع بوّابة رؤية \en{SEC-02}: لا اشتراك إلّا في
  فئة نشطة أو منشور مرئيّ فعلًا.
\item الحظر يُنشَأ بصفّ \code{UserBlock} واحد، لكنّه يُنفَّذ دفاعيًّا في كل مسار: فحص عند
  الكتابة (\code{SendMessage}، \code{CreateReply})، وفلتر عند القراءة (\code{NotBlockedWith})
  — وكلّه ثنائيّ الاتجاه فلا ثغرة.
\end{itemize}
\end{summary}
